%% Fichier contenant une annexe. La classe ne genere qu'un seul
%% fichier annexe.tex par defaut. Si vous en avez besoin davantage,
%% enregistrez ce fichier sous un autre nom et incluez-le dans votre
%% gabarit avec la commande \include.
%%
%% Pour que la bibliographie et la (les) annexe(s) soient paginees
%% correctement, c'est-a-dire en chiffres arabes pour la bibliographie
%% et en chiffres romains pour les annexes, ces dernieres doivent etre
%% placees apres la commande \backmatter. Ce faisant, la numerotation
%% des annexes s'en trouve desactivee. Vous devrez donc numeroter
%% vos annexes manuellement a l'interieur de la commande \chapter.
\chapter{Annexe A -- Titre de l'annexe}

%% Redigez votre annexe ici.

%% THESES ET MEMOIRES PAR ARTICLES SEULEMENT
%% Si vous avez insere des citations dans cette section, retirez les signes
%% de commentaires (%%) devant la commande ci-dessous et inscrivez le style
%% bibliographique et le nom du fichier .bib utilises pour vos references.
%% \HECreferences{style}{nom-du-fichier-file-name}
